% \iffalse meta-comment
% 
% This is file `bicaption.dtx'.
% 
% Copyright (C) 2010-2011 Axel Sommerfeldt (axel.sommerfeldt@f-m.fm)
% 
% --------------------------------------------------------------------------
% 
% This work may be distributed and/or modified under the
% conditions of the LaTeX Project Public License, either version 1.3
% of this license or (at your option) any later version.
% The latest version of this license is in
%   http://www.latex-project.org/lppl.txt
% and version 1.3 or later is part of all distributions of LaTeX
% version 2003/12/01 or later.
% 
% This work has the LPPL maintenance status "maintained".
% 
% This Current Maintainer of this work is Axel Sommerfeldt.
% 
% This work consists of the files caption.ins, caption.dtx, caption2.dtx,
% caption3.dtx, bicaption.dtx, ltcaption.dtx, and subcaption.dtx,
% the derived files caption.sty, caption2.sty, caption3.sty, bicaption.sty,
% ltcaption.sty, and subcaption.sty, and the user manuals caption-deu.tex,
% caption-eng.tex, and caption-rus.tex.
% 
% \fi
% \CheckSum{188}
%
% \iffalse
%<*driver>
\NeedsTeXFormat{LaTeX2e}[1994/12/01]
\ProvidesFile{bicaption.drv}[2010/09/19 v1.0 Adds a bilingual caption feature to the caption package]
\hbadness=9999 \newcount\hbadness \hfuzz=74pt % Make TeX shut up.
%\errorcontextlines=3
%
\documentclass[german,english]{ltxdoc}
\setlength\parindent{0pt}
\setlength\parskip{\smallskipamount}
%
\newcommand\LineBreak{\linebreak[3]}
\newcommand\PageBreak{\pagebreak[3]}
\usepackage{ifpdf}
\ifpdf
  \usepackage{mathptmx,courier}
  \usepackage[scaled=0.90]{helvet}
  \addtolength\marginparwidth{15pt}
  \ifdim\paperheight=297mm % a4paper
    \renewcommand\LineBreak{\\}
    \renewcommand\PageBreak{\clearpage}
  \fi
\fi
%
\usepackage[T1]{fontenc}
\usepackage[german,english]{babel}
\usepackage{selinput}\SelectInputMappings{adieresis={ä},germandbls={ß}}
%
\usepackage[bottom]{footmisc}
%
\usepackage{hypdoc}
\ifpdf\usepackage{hypdestopt}\fi
\hypersetup{pdfkeywords={LaTeX, package, bicaption},pdfstartpage={},pdfstartview={}}
%
\usepackage{bicaption}[2010/09/01]
\usepackage{subcaption}[2010/09/01]
%
\newcommand*\purerm[1]{{\upshape\mdseries\rmfamily #1}}
\newcommand*\puresf[1]{{\upshape\mdseries\sffamily #1}}
\newcommand*\purett[1]{{\upshape\mdseries\ttfamily #1}}
\let\package\puresf
\let\env\purett \let\opt\purett
%
\newcommand*\csmarg[1]{\texttt{\char`\{#1\char`\}}}
\newcommand*\csoarg[1]{\texttt{\char`\[#1\char`\]}}
\newcommand*\version[2][]{$v#2$}
\newcommand*\x{\discretionary{}{}{}}
%
\usepackage{marvosym}
\makeatletter
\newcommand*\INFO{\@ifstar{\@INFO{}}{\@INFO{\vbox to \ht\strutbox}}}
\newcommand*\@INFO[1]{\MARGINSYM{#1{\LARGE\Info}}}
\makeatother
%
\usepackage[alpine]{ifsym}
\newenvironment{background}{\par\bigskip\csname background*\endcsname}{\csname endbackground*\endcsname}
\newenvironment{background*}{\small\MARGINSYM{\Mountain}\ignorespaces}{\par}
%
\newcommand*\MARGINSYM[1]{\hskip 1sp \marginpar{\raggedleft\textcolor{blue}{{#1}}}}
%\newcommand*\NEW[2]{}%\hskip 1sp \marginpar{\footnotesize\sffamily\raggedleft#1\\#2}}
%
\newenvironment{Options}[1]%
  {\list{}{\renewcommand\makelabel[1]{\texttt{##1}\hfil}%
   \settowidth\labelwidth{\texttt{#1\space}}%
   \setlength\leftmargin{10pt}%
   \addtolength\leftmargin{\labelwidth}%
   \addtolength\leftmargin{\labelsep}}}%
  {\endlist}
%
\begin{document}
  \DocInput{bicaption.dtx}
\end{document}
%</driver>
% \fi
%
% \let\subsectionautorefname\sectionautorefname
% \let\subsubsectionautorefname\sectionautorefname
%
% \def\thispackage{the \package{bicaption} package}
% \def\Thispackage{The \package{bicaption} package}
%
% \newcommand\NEWfeature{\NEW{New feature}}
% \newcommand\NEWdescription{\NEW{New description}}
%
% \makeatletter
% \newcommand*\Ref{\@ifstar{\@Ref\ref}{\@Ref\autoref}}
% \newcommand*\@Ref[2]{#1{#2}: \textit{\nameref{#2}}}
% \makeatother
% \newcommand*\See[1]{\nopagebreak{\small (See #1)}}
%
% \GetFileInfo{bicaption.drv}
% \let\docdate\filedate
% \GetFileInfo{bicaption.sty}
%
% \title{\texorpdfstring{\Thispackage\thanks{%^^A
%          This package has version number \fileversion, last revised \filedate.}}%^^A
%        {The bicaption package}}
% \author{Axel Sommerfeldt\\
%         \href{mailto:axel.sommerfeldt@f-m.fm}{\texttt{axel.sommerfeldt@f-m.fm}}}
% \date{\docdate}
% \maketitle
% 
% \begin{abstract}
% This package supports the typesetting of bilangual captions.
% \end{abstract}
%
% \setcounter{tocdepth}{2}
% \tableofcontents
%
% \clearpage
% \section{The user interface}
%
% \subsection{Loading the package}
%
% \DescribeMacro{\usepackage}
% This package will be loaded by
% \begin{quote}
%   |\usepackage|\oarg{options}|{bicaption}|\quad.
% \end{quote}
% The options for \thispackage\ are the same ones as for the \package{caption}
% package and specify settings which are used for the second language
% \emph{additionally}.
% In fact
% \begin{quote}
%   |\usepackage|\oarg{options}|{bicaption}|
% \end{quote}
% is identical to
% \begin{quote}
%   |\usepackage{bicaption}|\\
%   |\captionsetup[bi-second]|\marg{options}\quad.
% \end{quote}
%
% \subsection{Setting options}
%
% \DescribeMacro{\captionsetup}
% \begin{quote}
%   |\captionsetup[bi]|\marg{options}
% \end{quote}
% do setup options which will be used for bilanguage captions \emph{additionally}
% to the ones which are setup for the specific floating environment.
%
% \begin{quote}
%   |\captionsetup[bi-first]|\marg{options}
% \end{quote}
% do setup options which will be used for the \emph{first} heading
% of the bilanguage captions \emph{additionally}
% to the ones which are setup for the specific floating environment
% and the ones which are setup by |\captionsetup[bi]{|\ldots|}|.
%
% \begin{quote}
%   |\captionsetup[bi-second]|\marg{options}
% \end{quote}
% do setup options which will be used for the \emph{second} heading
% of the bilanguage captions \emph{additionally}
% to the ones which are setup for the specific floating environment
% and the ones which are setup by |\captionsetup[bi]{|\ldots|}|.
%
% \bigskip
%
% Options specified with |\usepackage[|\ldots|]{bi|\-|caption}| and
% |\caption|\-|setup[bi|\ldots|]{|\ldots|}| will override the ones specified by
% |\caption|\-|setup{|\ldots|}| and |\caption|\-|setup[fig|\-|ure]{|\ldots|}|
% (same for `table'). So finally we have the following order how
% settings for bilingual captions are applied:
% \begin{enumerate}
% \item Global settings
%  {\small(|\usepackage[|\ldots|]{caption}| and |\captionsetup{|\ldots|}|)}
% \item Environmental settings
%  {\small(|\captionsetup[figure|\emph{ -or- }|table]{|\ldots|}|)}
% \item Local settings
%  {\small(|\captionsetup{|\ldots|}| inside |figure| or |table| environment)}
% \item Custom `bi' settings
%  {\small(|\captionsetup[bi]{|\ldots|}|)}
% \item Custom `bi-first' resp.~`bi-second' settings
%  {\small(|\usepackage[|\ldots|]{bicaption}| and
%          |\caption|\-|setup[bi-first]{|\ldots|}|
%    resp.~|\caption|\-|setup[bi-second]{|\ldots|}|)}
% \end{enumerate}
% An example:
% \begin{quote}
%   |\usepackage[labelsep=quad,indention=10pt]{caption}|\\
%   |\usepackage[labelfont=bf]{bicaption}|\\
%   |\captionsetup[table]{labelfont=it,position=top}|
% \end{quote}
% causes the second heading of the bilingual caption inside |table| environments
% to be typeset with the settings
% \begin{quote}
%   |labelsep=quad,indention=10pt,position=top,labelfont=bf|~.
% \end{quote}
%
% \subsection{Additional options}
%
% These options are available additional to the ones offered by the
% \package{caption} package:
%
% \begin{Options}{language=}
%   \item[language=]
%   \DescribeMacro{language=}
%   Sets the language of the caption,
%   e.g.
%   \begin{quote}|\use|\-|package|\-|[eng|\-|lish]{bi|\-|caption}|\end{quote}
%   will typeset the second caption of bilingual captions in English.
%   (The language will be set with |\select|\-|language| internally,
%    so the \package{babel} package must be loaded for using this option.)
%
%   \item[bi-lang=]
%   \DescribeMacro{bi-lang=}
%   Causes a selection of the headings of bilingual captions.
%   \begin{quote}|\captionsetup{bi-lang=both}|\end{quote}
%   will cause that both caption headings are being typeset.\\
%   (This is the default.)
%   \begin{quote}|\captionsetup{bi-lang=first}|\end{quote}
%   will cause that only the \emph{first} heading is being typeset, and
%   \begin{quote}|\captionsetup{bi-lang=second}|\end{quote}
%   will cause that only the \emph{second} heading is being typeset.
%
%   \item[bi-slc=]
%   \DescribeMacro{bi-singlelinecheck=}
%   Switches the common single-line-check |on| or |off|, i.e.~when switched on
%   only a single check will be done for both captions, and the result will affect
%   both captions afterwards. So if only one caption is longer than a single line,
%   both captions will be treated as if they are longer than a single line, even if
%   the second one isn't. (The default is |on|.)
%
%   \item[bi-swap=]
%   \DescribeMacro{bi-swap=}
%   \begin{quote}|\captionsetup{bi-swap}|\end{quote}
%   will swap the primary and secondary language,
%   making the first language the second one and vice versa. (The default is |false|.)
% \end{Options}
%
% \subsection{The \cs{bicaption} commands}
%
% \DescribeMacro\bicaption
% Bilingual captions will be typeset by
% \begin{quote}
%   |\bicaption|\oarg{list entry}\marg{heading \#1}\marg{heading \#2}\\
%   |\bicaption*|\marg{heading \#1}\marg{heading \#2}
% \end{quote}
% The |\label| should be placed either after this command, or inside the first heading.
%
% \bigskip
%
% If the \package{subcaption} package is loaded, these commands are available
% additionally:
%
% \DescribeMacro\bisubcaption
% Bilingual sub-captions will be typeset by
% \begin{quote}
%   |\bisubcaption|\oarg{list entry}\marg{heading \#1}\marg{heading \#2}\\
%   |\bisubcaption*|\marg{heading \#1}\marg{heading \#2}
% \end{quote}
% The |\label| should be placed either after this command, or inside the first heading.
%
% \DescribeMacro\bisubcaptionbox
% Bilingual sub-caption boxes will be typeset by
% \begin{quote}
%   |\subcaptionbox|\oarg{list entry}\marg{heading \#1}\marg{heading \#2}\\
%   |              |\oarg{width}\oarg{inner-pos}\marg{contents}\\
%   |\subcaptionbox*|\marg{heading \#1}\marg{heading \#2}\\
%   |               |\oarg{width}\oarg{inner-pos}\marg{contents}
% \end{quote}
% The |\label| should be placed inside the first heading.
%
% \subsection{An example code}
%
% \begin{quote}
%   |\documentclass[english,ngerman]{article}|\\
%   |\usepackage{selinput}|\\
%   |\SelectInputMappings{adieresis={ä},germandbls={ß}}|\\
%   ||\\
%   |\usepackage{babel}|\\
%   |\usepackage[lang=english,font=it]{bicaption}|\\
%   |\usepackage[format=hang,list=on]{subcaption}|\\
%   ||\\
%   |\begin{document}|\\
%   ||\\
%   \iffalse
%   |\listoffigures|\\
%   ||\\
%   \fi
%   |\begin{figure}[!htb]|\\
%   |  \centering|\\
%   |  \bisubcaptionbox|\\
%   |    {Teilabbildung A\label{fig:test:A}}|\\
%   |    {Subfigure A}[0.4\textwidth]{IMAGE}%|\\
%   |  \qquad|\\
%   |  \bisubcaptionbox|\\
%   |    {Teilabbildung langer Titel B\label{fig:test:B}}|\\
%   |    {Subfigure long title B}[0.4\textwidth]{IMAGE}%|\\
%   |  \bicaption{Deutscher Titel}{English Title}|\\
%   |  \label{fig:test}|\\
%   |\end{figure}|\\
%   ||\\
%   \iffalse
%   |\captionsetup{bi-lang=2nd}|\\
%   ||\\
%   |\begin{figure}[!htb]|\\
%   |  \centering|\\
%   |  \bisubcaptionbox|\\
%   |    {Teilabbildung A\label{fig:test2:A}}|\\
%   |    {Subfigure A}[0.4\textwidth]{IMAGE}%|\\
%   |  \qquad|\\
%   |  \bisubcaptionbox|\\
%   |    {Teilabbildung langer Titel B\label{fig:test2:B}}|\\
%   |    {Subfigure long title B}[0.4\textwidth]{IMAGE}%|\\
%   |  \bicaption{Deutscher Titel}{English Title}|\\
%   |  \label{fig:test2}|\\
%   |\end{figure}|\\
%   \fi
%   ||\\
%   |\captionsetup{bi-lang=both}|\\
%   ||\\
%   |\begin{figure}[!htb]|\\
%   |  \centering|\\
%   |  \bisubcaptionbox[A]|\\
%   |    {Und eine gaaaanz lange Caption: Teilabbildung A}|\\
%   |    {Subfigure A}[0.4\textwidth]{IMAGE}%|\\
%   |  \qquad|\\
%   |  \bisubcaptionbox[B]|\\
%   |    {Teilabbildung B}|\\
%   |    {Subfigure B}[0.4\textwidth]{IMAGE}%|\\
%   |  \bicaption[Abbildungsverzeichnistitel]|\\
%   |    {Und eine noch viel viel viel|\\
%   |     längere deutsche Beschriftung: Deutscher Titel}|\\
%   |    {Short English heading}|\\
%   |\end{figure}|\\
%   ||\\
%   |\captionsetup{bi-slc=0}|\\
%   ||\\
%   |\begin{figure}[!htb]|\\
%   |  \centering|\\
%   |  \bisubcaptionbox[A]|\\
%   |    {Und eine gaaaanz lange Caption: Teilabbildung A}|\\
%   |    {Subfigure A}[0.4\textwidth]{IMAGE}%|\\
%   |  \qquad|\\
%   |  \bisubcaptionbox[B]|\\
%   |    {Teilabbildung B}|\\
%   |    {Subfigure B}[0.4\textwidth]{IMAGE}%|\\
%   |  \bicaption[Abbildungsverzeichnistitel]|\\
%   |    {Und eine noch viel viel viel|\\
%   |     längere deutsche Beschriftung: Deutscher Titel}|\\
%   |    {Short English heading}|\\
%   |\end{figure}|\\
%   ||\\
%   |\captionsetup{slc=0}|\\
%   ||\\
%   |\begin{figure}[!htb]|\\
%   |  \centering|\\
%   |  \bisubcaptionbox[A]|\\
%   |    {Und eine gaaaanz lange Caption: Teilabbildung A}|\\
%   |    {Subfigure A}[0.4\textwidth]{IMAGE}%|\\
%   |  \qquad|\\
%   |  \bisubcaptionbox[B]|\\
%   |    {Teilabbildung B}|\\
%   |    {Subfigure B}[0.4\textwidth]{IMAGE}%|\\
%   |  \bicaption[Abbildungsverzeichnistitel]|\\
%   |    {Und eine noch viel viel viel|\\
%   |     längere deutsche Beschriftung: Deutscher Titel}|\\
%   |    {Short English heading}|\\
%   |\end{figure}|\\
%   ||\\
%   |\end{document}|\\
% \end{quote}
%
% \captionsetup[bi-first]{lang=german}
% \captionsetup[bi-second]{lang=english,font=it}
% \captionsetup[sub]{format=hang,list=on}
%
% \iffalse
% \listoffigures
% \bigskip
% \fi
%
% \begin{figure}[!htb]
%   \centering
%   \bisubcaptionbox
%     {Teilabbildung A\label{fig:test:A}}
%     {Subfigure A}[0.4\textwidth]{IMAGE}%
%   \qquad
%   \bisubcaptionbox
%     {Teilabbildung langer Titel B\label{fig:test:B}}
%     {Subfigure long title B}[0.4\textwidth]{IMAGE}%
%   \bicaption{Deutscher Titel}{English Title}
%   \label{fig:test}
% \end{figure}
%
% \iffalse
% \captionsetup{bi-lang=2nd}
%
% \begin{figure}[!htb]
%   \centering
%   \bisubcaptionbox
%     {Teilabbildung A\label{fig:test2:A}}
%     {Subfigure A}[0.4\textwidth]{IMAGE}%
%   \qquad
%   \bisubcaptionbox
%     {Teilabbildung langer Titel B\label{fig:test2:B}}
%     {Subfigure long title B}[0.4\textwidth]{IMAGE}%
%   \bicaption{Deutscher Titel}{English Title}
%   \label{fig:test2}
% \end{figure}
% \fi
%
% \captionsetup{bi-lang=both}
%
% \begin{figure}[!htb]
%   \centering
%   \bisubcaptionbox[A]
%     {Und eine gaaaanz lange Caption: Teilabbildung A}
%     {Subfigure A}[0.4\textwidth]{IMAGE}%
%   \qquad
%   \bisubcaptionbox[B]
%     {Teilabbildung B}
%     {Subfigure B}[0.4\textwidth]{IMAGE}%
%   \bicaption[Abbildungsverzeichnistitel]
%     {Und eine noch viel viel viel
%      längere deutsche Beschriftung: Deutscher Titel}
%     {Short English heading}
% \end{figure}
%
% \captionsetup{bi-slc=0}
%
% \begin{figure}[!htb]
%   \centering
%   \bisubcaptionbox[A]
%     {Und eine gaaaanz lange Caption: Teilabbildung A}
%     {Subfigure A}[0.4\textwidth]{IMAGE}%
%   \qquad
%   \bisubcaptionbox[B]
%     {Teilabbildung B}
%     {Subfigure B}[0.4\textwidth]{IMAGE}%
%   \bicaption[Abbildungsverzeichnistitel]
%     {Und eine noch viel viel viel
%      längere deutsche Beschriftung: Deutscher Titel}
%     {Short English heading}
% \end{figure}
%
% \captionsetup{slc=0}
%
% \begin{figure}[!htb]
%   \centering
%   \bisubcaptionbox[A]
%     {Und eine gaaaanz lange Caption: Teilabbildung A}
%     {Subfigure A}[0.4\textwidth]{IMAGE}%
%   \qquad
%   \bisubcaptionbox[B]
%     {Teilabbildung B}
%     {Subfigure B}[0.4\textwidth]{IMAGE}%
%   \bicaption[Abbildungsverzeichnistitel]
%     {Und eine noch viel viel viel
%      längere deutsche Beschriftung: Deutscher Titel}
%     {Short English heading}
% \end{figure}
%
% \iffalse
% --------------------------------------------------------------------------- %
% \fi
%
% \StopEventually{%^^A
% }
%
% \iffalse
% --------------------------------------------------------------------------- %
% \fi
%
% \DoNotIndex{\\,\_,\ ,\@@par}
% \DoNotIndex{\@bsphack}
% \DoNotIndex{\@car,\@cdr,\@classoptionslist,\@cons,\@currext,\@currname}
% \DoNotIndex{\@ehc,\@ehd,\@empty,\@esphack,\@expandtwoargs}
% \DoNotIndex{\@for,\@firstofone,\@firstoftwo}
% \DoNotIndex{\@gobble,\@gobblefour,\@gobbletwo,\@hangfrom}
% \DoNotIndex{\@ifnextchar,\@ifpackagelater,\@ifpackageloaded}
% \DoNotIndex{\@ifstar,\@ifundefined,\@latex@error,\@namedef,\@nameuse}
% \DoNotIndex{\@onlypreamble,\@parboxrestore,\@plus,\@ptionlist}
% \DoNotIndex{\@removeelement,\@restorepar,\@secondoftwo,\@setpar}
% \DoNotIndex{\@tempa,\@tempboxa,\@tempdima,\@tempdimb,\@tempdimc,\@tempb,\@tempc}
% \DoNotIndex{\@testopt}
% \DoNotIndex{\@undefined,\@unprocessedoptions,\@unusedoptionlist}
% \DoNotIndex{\p@,\z@}
% \DoNotIndex{\active,\addtocounter,\addtolength,\advance,\aftergroup}
% \DoNotIndex{\baselineskip,\begin,\begingroup,\bfseries,\box}
% \DoNotIndex{\catcode,\centering,\changes,\csname,\def,\divide,\do,\downarrow}
% \DoNotIndex{\edef,\else,\empty,\end,\endcsname,\endgraf,\endgroup,\expandafter}
% \DoNotIndex{\fi,\footnotesize,\global}
% \DoNotIndex{\hangindent,\hbox,\hfil,\hsize,\hskip,\hspace,\hss}
% \DoNotIndex{\ifcase,\ifdim,\ifnum,\ifodd,\ifvoid,\ifvmode}
% \DoNotIndex{\ifx,\ignorespaces,\itshape}
% \DoNotIndex{\Large,\large,\leavevmode,\leftmargini,\leftskip,\let,\linewidth}
% \DoNotIndex{\llap,\long,\m@ne,\margin,\mdseries,\message}
% \DoNotIndex{\newcommand,\newdimen,\newlength,\newline,\newif,\newsavebox}
% \DoNotIndex{\next,\nobreak,\nobreakspace,\noexpand,\noindent,\numberline}
% \DoNotIndex{\normalcolor,\normalfont,\normalsize,\or,\par,\parbox,\parfillskip}
% \DoNotIndex{\parindent,\parskip,\prevdepth,\protect,\protected@edef,\protected@write}
% \DoNotIndex{\providecommand,\quad}
% \DoNotIndex{\raggedleft,\raggedright,\relax,\renewcommand,\RequirePackage}
% \DoNotIndex{\rightskip,\rmfamily}
% \DoNotIndex{\sbox,\scriptsize,\scshape,\setbox,\setlength,\sffamily,\slshape}
% \DoNotIndex{\small,\string,\space,\strut}
% \DoNotIndex{\textheight,\the,\toks@,\typeout,\ttfamily}
% \DoNotIndex{\unvbox,\uparrow,\upshape,\usebox,\usepackage}
% \DoNotIndex{\value,\vbox,\vsize,\vskip,\wd,\width,\z@skip}
% \DoNotIndex{\AtBeginDocument,\AtEndOfPackage,\CurrentOption,\DeclareOption}
% \DoNotIndex{\ExecuteOptions,\GenericWarning,\IfFileExists,\InputIfFileExists}
% \DoNotIndex{\NeedsTeXFormat,\MessageBreak}
% \DoNotIndex{\PackageError,\PackageInfo,\PackageWarning,\PackageWarningNoLine}
% \DoNotIndex{\PassOptionsToPackage,\ProcessOptions,\ProvidesPackage}
%
% \iffalse
% --------------------------------------------------------------------------- %
% \fi
%
% \setlength{\parskip}{0pt plus 1pt}
% \changes{v0.1}{2010/07/13}{Initial version}
% \changes{v0.2}{2010/07/13}{Check for caption package added}
% \changes{v0.3}{2010/07/13}{Usage of \cs{caption@applyfont} added}
% \changes{v0.4}{2010/07/13}{``Singlelinecheck'' fixed}
% \changes{v0.5}{2010/07/13}{Options \opt{bi-first} and \opt{bi-second} added}
% \changes{v0.6}{2010/07/13}{Option \opt{bi-slc} added}
% \changes{v0.7}{2010/07/13}{Option \opt{bi-lang} added}
% \changes{v0.8}{2010/09/04}{Adapted to current version of the caption kernel}
% \changes{v0.9}{2010/09/17}{Option \opt{bi-swap} added}
% \changes{v0.9a}{2011/07/13}{Warning regarding \package{babel} package added}
%
% \newcommand*\Note[2][Note]{\par{\small\emph{#1:} #2}}
%
% \iffalse
% --------------------------------------------------------------------------- %
% \fi
%
% \clearpage
% \section{The implementation}
% \iffalse
%<*package>
% \fi
%
% \subsection{Identification}
%
%    \begin{macrocode}
\NeedsTeXFormat{LaTeX2e}[1994/12/01]
\ProvidesPackage{bicaption}[2011/07/13 v0.9a Bilingual Captions (AR)]
\RequirePackage{caption}[2011/01/01] % needs v3.2 or newer
%    \end{macrocode}
% \bigskip
%
% \pagebreak[3]
% \subsection{Initial code}
%
% \begin{macro}{\bicaption@Warning}
%  |\bicaption@Warning|\marg{message}
%    \begin{macrocode}
\newcommand*\bicaption@Warning[1]{%
  \bicaption@WarningNoLine{#1\on@line}}
%    \end{macrocode}
% \end{macro}
% \begin{macro}{\bicaption@WarningNoLine}
%  |\bicaption@WarningNoLine|\marg{message}
%    \begin{macrocode}
\newcommand*\bicaption@WarningNoLine[1]{%
  \PackageWarning{bicaption}{#1.^^J\bicaption@wh\@gobbletwo}}
%    \end{macrocode}
%    \begin{macrocode}
\newcommand*\bicaption@wh{%
  See the bicaption package documentation for explanation.}
%    \end{macrocode}
% \end{macro}
% \begin{macro}{\bicaption@Error}
%  |\bicaption@Error|\marg{message}
%    \begin{macrocode}
\newcommand*\bicaption@Error[1]{%
  \PackageError{bicaption}{#1}\bicaption@eh}
%\let\bicaption@KV@err\bicaption@Error
%    \end{macrocode}
%    \begin{macrocode}
\newcommand*\bicaption@eh{%
  If you do not understand this error, please take a closer look\MessageBreak
  at the documentation of the `bicaption' package.\MessageBreak\@ehc}
%    \end{macrocode}
% \end{macro}
%
% \pagebreak[3]
% \subsection{Declaration of options}
%
% The option |bi-lang| will setup which language(s) will actually be typeset,
% the first one, the second one, or both of them.
%    \begin{macrocode}
\newcount\bicaption@lang
\DeclareCaptionOption{bi-lang}{%
  \caption@ifinlist{#1}{0,all,both}{%
    \bicaption@lang=0\relax
  }{\caption@ifinlist{#1}{1,1st,first}{%
    \bicaption@lang=1\relax
  }{\caption@ifinlist{#1}{2,2nd,second}{%
    \bicaption@lang=2\relax
  }{%
    \bicaption@Error{Undefined bi-lang value `#1'}%
  }}}}
%    \end{macrocode}
%
% The option |bi-singlelinecheck| will setup if a single check will be used
% for both languages (|=on|),
% or if both languages will be checked individually (|=off|).
%    \begin{macrocode}
\DeclareCaptionOption{bi-singlelinecheck}[1]{%
  \caption@set@bool\bicaption@ifslc{#1}}
\DeclareCaptionOption{bi-slc}[1]{%
  \caption@set@bool\bicaption@ifslc{#1}}
%    \end{macrocode}
%
% The option |bi-swap| will swap the primary and secondary language,
% making the first language the second one and vice versa.
%    \begin{macrocode}
\DeclareCaptionOption{bi-swap}[1]{%
  \caption@set@bool\bicaption@ifswap{#1}}
%    \end{macrocode}
%
% The option |lang=|\meta{language} will setup the language of the caption.
%    \begin{macrocode}
\DeclareCaptionOption{lang}{\def\bicaption@language{#1}}
\let\KV@caption@language\KV@caption@lang
%    \end{macrocode}
%
% \begin{macro}{\bicaption@selectlanguage}
% Set the language via \cs{selectlanguage}.
%    \begin{macrocode}
\newcommand*\bicaption@selectlanguage{%
  \@ifundefined{bicaption@language}{}{%
    \expandafter\selectlanguage\expandafter{\bicaption@language}}}
%    \end{macrocode}
% |\caption@applyfont| (of the \package{caption} package kernel) will be extended here
% so the language setting will actually take effect.
%    \begin{macrocode}
\g@addto@macro\caption@applyfont{%
  \bicaption@selectlanguage}
%    \end{macrocode}
%    \begin{macrocode}
\g@addto@macro\caption@prepareslc{%
  \let\bicaption@language\@undefined}
%    \end{macrocode}
% \end{macro}
%
% \pagebreak[3]
% \subsection{Execution of options}
%
% Setup default values for |bi-lang| and |bi-singlelinecheck|.
%    \begin{macrocode}
\caption@ExecuteOptions{caption}{bi-lang=0,bi-slc=1,bi-swap=0}
%    \end{macrocode}
%
% Set the language for the first caption.
%    \begin{macrocode}
\ifx\bbl@main@language\@undefined
  \bicaption@WarningNoLine{Please load this package after the babel package}
\else
  \edef\@tempa{\noexpand\captionsetup[bi-first]{lang=\bbl@main@language}}
  \@tempa
\fi
%    \end{macrocode}
% We use |\caption@ProcessOptions| here to add the options to the `|bi-second|' option
% list instead of executing them immediately.
%    \begin{macrocode}
\caption@SetupOptions{bicaption}{\captionsetup[bi-second]{#2}}%
\caption@ProcessOptions*{bicaption}
%    \end{macrocode}
%
% \pagebreak[3]
% \subsection{Main code}
%
% \begin{macro}{\caption@@make}
% We patch \cs{caption@@make} (of the \package{caption} package kernel)
% so \cs{bicaption@@make} will be used for bilingual captions instead.
%    \begin{macrocode}
\let\caption@@make@ORI\caption@@make
\renewcommand\caption@@make[2]{%
  \@ifundefined{bicaption@text}%
    {\caption@@make@ORI{#1}{#2}}%
    {\bicaption@@make{#1}{#2}{\bicaption@text}%
     \global\let\bicaption@label\@undefined
     \global\let\bicaption@text\@undefined}}
%    \end{macrocode}
% \end{macro}
%
% \begin{macro}{\bicaption@@make}
% |\bicaption@@make|\marg{label}\marg{text \#1}\marg{text \#2}\par
% Typeset both captions using the original version of \cs{caption@@make}.
%    \begin{macrocode}
\newcommand\bicaption@@make[3]{%
%    \end{macrocode}
% Execute the options setup with |\captionsetup[bi]{|\ldots|}|.
%    \begin{macrocode}
  \caption@setoptions{bi}%
%    \end{macrocode}
% Perform the common single-line-check for both captions, if requested.
%    \begin{macrocode}
  \ifnum\bicaption@lang=0\relax
    \bicaption@ifslc
      {\caption@@slc{#1}{#2}{\captionwidth}{}%
                    {\caption@set@bool\caption@ifslc0}%
       \caption@@slc{#1}{#3}{\captionwidth}{}%
                    {\caption@set@bool\caption@ifslc0}}%
      {}%
  \fi
%    \end{macrocode}
% Typeset the first caption, if requested.
% (Otherwise we only apply the label of it.)
%    \begin{macrocode}
  \ifnum\bicaption@lang=2\relax
    \ifx\bicaption@label\@empty\else
      \expandafter\label\expandafter{\bicaption@label}%
    \fi
  \else
    \begingroup
      \caption@setoptions{bi-first}%
      \caption@@make@ORI{#1}{#2}%
    \endgroup
  \fi
%    \end{macrocode}
% Typeset the second caption, if requested.
%    \begin{macrocode}
  \ifnum\bicaption@lang=1\relax
  \else
    \begingroup
      \caption@setoptions{bi-second}%
      \caption@@make@ORI{#1}{#3}%
    \endgroup
  \fi}
%    \end{macrocode}
% \end{macro}
%
% \begin{macro}{\bicaption@setup}
% |\bicaption@setup|\marg{text \#1}\marg{text \#2}\par
% Initiates the bilingual caption typesetting by extracting the |\label|
% of the first text and storing the second text into |\bicaption@text|.
%    \begin{macrocode}
\newcommand\bicaption@setup[2]{%
  \bicaption@getlabel#1\label{}\@nil
  \global\long\def\bicaption@text{\ignorespaces#2}}
\long\def\bicaption@getlabel#1\label#2#3\@nil{%
  \global\def\bicaption@label{#2}}
%    \end{macrocode}
% \end{macro}
%
% \pagebreak[3]
% \subsubsection{The \cs{bicaption} commands}
%
% \begin{macro}{\@bicaption}
%    \begin{macrocode}
\newcommand*\@bicaption[1]{%
  \let\bicaption@cmd#1%
  \caption@withoptargs\@@bicaption}
\newcommand\@@bicaption[3]{%
  \bicaption@ifswap
    {\bicaption@setup{#2}{#2}%
     \bicaption@cmd#1{#3}}%
    {\bicaption@setup{#2}{#3}%
     \bicaption@cmd#1{#2}}}
%    \end{macrocode}
% \end{macro}
%
% \begin{macro}{\bicaption}
% |\bicaption*|\oarg{list entry}\marg{text \#1}\marg{text \#2}
%    \begin{macrocode}
\newcommand\bicaption{\@bicaption\caption}
%    \end{macrocode}
% \end{macro}
%
% \begin{macro}{\bisubcaption}
% |\bisubcaption*|\oarg{list entry}\marg{text \#1}\marg{text \#2}
%    \begin{macrocode}
\newcommand\bisubcaption{\@bicaption\subcaption}
%    \end{macrocode}
% \end{macro}
%
% \begin{macro}{\bisubcaptionbox}
% |\bisubcaptionbox*|\oarg{list entry}\marg{text \#1}\marg{text \#2}%
%                    \oarg{\ldots}\marg{\ldots}
%    \begin{macrocode}
\newcommand\bisubcaptionbox{\@bicaption\subcaptionbox}
%    \end{macrocode}
% \end{macro}
%
% \iffalse
%</package>
% \fi
%
% \iffalse
% --------------------------------------------------------------------------- %
% \fi
%
% \Finale
%
\endinput
